\subsection{Assignment 10 -- Confronto regionale delle performance di streaming}

L’Assignment 10 analizza il rapporto tra gli stream di una categoria musicale
in una determinata regione e gli stream della stessa categoria nelle restanti
regioni, considerando esclusivamente le canzoni in cui l’artista ricopre il ruolo
di \textit{primary artist}.

La logica analitica è stata inizialmente validata tramite una formulazione SQL
al fine di verificarne la correttezza semantica e numerica, e successivamente
implementata in \textit{SQL Server Integration Services} (SSIS) mediante un unico
\textit{Data Flow}. Tale scelta garantisce coerenza con lo schema del
\textit{Data Warehouse} e consente di mantenere un flusso di trasformazioni
efficiente e facilmente tracciabile.

Il flusso dati ha origine dalla fact table \textit{FactParticipation} ed è
progressivamente arricchito tramite operazioni di \textit{Lookup} sulle
dimensioni \textit{DimSong}, \textit{DimArtist} e \textit{DimArtistGeography}.
Queste trasformazioni permettono di associare a ciascun record la categoria
musicale del brano e la regione di nascita dell’artista principale. Un
\textit{Conditional Split} viene utilizzato per selezionare esclusivamente le
righe relative ai \textit{primary artist}, in accordo con la consegna.

L’aggregazione degli stream per ciascuna coppia (categoria, regione) è realizzata
tramite una trasformazione di \textit{Aggregate}. In parallelo, mediante una
\textit{Multicast}, viene calcolato il totale degli stream per categoria, valore
necessario per determinare il volume di stream della stessa categoria nelle
altre regioni.

I risultati intermedi vengono successivamente ricombinati attraverso una
trasformazione di \textit{Merge Join}, che consente di calcolare sia gli stream
nelle restanti regioni sia il rapporto finale tra stream nella regione e stream
nelle altre regioni. Tali calcoli sono implementati mediante trasformazioni di
\textit{Derived Column}, con una gestione esplicita dei casi di denominatore nullo.

Prima della scrittura dell’output finale, è stata introdotta una trasformazione
di \textit{Data Conversion} sulle colonne \textit{Category} e \textit{Region}.
Questo passaggio è necessario poiché tali attributi provengono dalle dimensioni
con tipo Unicode (\textit{DT\_WSTR}), mentre la \textit{Flat File Destination}
richiede stringhe non Unicode (\textit{DT\_STR}). La conversione garantisce la
compatibilità dei tipi di dato e previene errori in fase di esecuzione, senza
alterare il contenuto informativo.
